% Generado por sistema/backend ReporteTex.java -- no editar a mano.
% Datos: CSV importado (titulados) / CSV importado (empleadores)
\subsubsection*{F3. Nivel del programa cursado}

Verificar si los programas de formación complementaria cursados se reparten por igual entre los niveles o si hay niveles más cursados.

\textit{Paso 1. Contar con la muestra y determinar su tamaño.} Población: titulados con formación complementaria; tamaño $n = 98$; opciones de la pregunta $k = 5$; $\alpha = 0.05$.

\begin{itemize}[nosep]
  \item $H_0$: Las respuestas de «Nivel del programa» se reparten por igual entre las 5 opciones (uniformidad).
  \item $H_1$: Las respuestas de «Nivel del programa» no se reparten por igual: hay opciones preferidas.
\end{itemize}

\textit{Paso 2. Calcular la frecuencia esperada.} Si las respuestas se repartieran por igual, cada opción tendría:
\[
FE = \frac{n}{k} = \frac{98}{5} = 19.60
\]

\textit{Paso 3. Calcular el estadístico chi-cuadrado.} Se cuenta la frecuencia observada $FO_i$ de cada opción y se calcula $\frac{(FO_i - FE)^2}{FE}$.

\begin{mitabla}{F3: prueba de las frecuencias, Nivel del programa}
\begin{tabular}{L{6cm}ccc}
\toprule
Opción & $FO_i$ & $FE$ & $\frac{(FO_i-FE)^2}{FE}$ \\
\midrule
Diplomado & 66 & 19.60 & $\frac{(66-19.60)^2}{19.60} = 109.8449$ \\[3pt]
Especialidad & 3 & 19.60 & $\frac{(3-19.60)^2}{19.60} = 14.0592$ \\[3pt]
Maestría & 27 & 19.60 & $\frac{(27-19.60)^2}{19.60} = 2.7939$ \\[3pt]
Doctorado & 2 & 19.60 & $\frac{(2-19.60)^2}{19.60} = 15.8041$ \\[3pt]
Pos doctorado & 0 & 19.60 & $\frac{(0-19.60)^2}{19.60} = 19.6000$ \\[3pt]
\midrule
Total & 98 & 98 & $x_0^2 = 162.1020$ \\
\bottomrule
\end{tabular}
\end{mitabla}

\[
x_0^2 = \sum_{i=1}^{5} \frac{(FO_i - FE)^2}{FE} = 162.1020
\]

\textit{Paso 4. Determinar el estadístico teórico.} $x^2_{\alpha,\,k-1} = x^2_{0.05,\,5-1} = x^2_{0.05,\,4} = 9.488$ (tabla chi-cuadrada).

\textit{Paso 5. Establecer la relación entre estadísticos.} $x_0^2 < x^2_{\alpha,\,k-1}$ $\rightarrow$ $162.1020 \ge 9.488$ $\rightarrow$ no.

\textit{Paso 6. Conclusión.}

\begin{center}
\begin{tabular}{l}
Si ($x_0^2 < x^2_{\alpha,\,k-1}$)\\
\quad ``Uniformidad''\\
Sino\\
\quad \fbox{No ``Uniformidad''}\\
FinSi
\end{tabular}
\end{center}

Se rechaza $H_0$: las respuestas no se reparten por igual, hay opciones preferidas. La más elegida es «Diplomado» (66 de 98, 67.3\%).

\textit{Recomendación.} Hay niveles más cursados: la oferta de actualización debe reforzar esos niveles, que son los que los titulados realmente siguen.

